% GENERATED FILE. Edit canonical lessons, then run npm run generate:books.
% canonical-source-hash: fnv1a64:9b222ef4cd675820
% canonical-lessons: ML-C151-kutira, ML-C151-kaluta, ML-C151-pottu, ML-C151-panni, ML-C151-simham

\chapter{Horse, Donkey, Buffalo, Pig, Lion}
\label{ch:ml-horse-donkey-buffalo-pig-lion}
\hlchaptermodality{151}

\begin{chapteropening}
\textbf{By the end of this chapter:} \emph{I can say a horse, a donkey, a buffalo, a pig and a lion.}

Complete the last lesson of chapter 151: I can say a horse, a donkey, a buffalo, a pig and a lion.
\end{chapteropening}

\section[\texorpdfstring{\ml{കുതിര} (kutira)}{കുതിര (kutira)}]{\ml{കുതിര} (kutira) — a horse}

\label{lesson:ML-C151-kutira}



\begin{quote}
Before the new one: say the Malayalam for black pepper, then the Malayalam for garlic.
\end{quote}

\subsection*{You'll want to know: \ml{കുതിര}}
\textbf{\ml{കുതിര}} — \emph{kutira} — \textquotedblleft{}a horse\textquotedblright{}.

\begin{grammarlens}[title={What you've built}]
One more word for this chapter.
\end{grammarlens}

\subsection*{Your turn}
\begin{itemize}
\raggedright
  \item \emph{Say it:} \emph{kutira}
  \item \emph{Say it:} \emph{kutira}, once more
  \item \emph{Say it:} say \emph{kutira}
  \item \emph{Recall:} say the Malayalam for black pepper, then the Malayalam for garlic, then say \emph{kutira} again
\end{itemize}

\begin{tcolorbox}[breakable,colback=teal!4,colframe=teal!35!black,title={Before you move on}]
What does \ml{കുതിര} mean? (\textquotedblleft{}A horse\textquotedblright{}.) Say it once more.
\end{tcolorbox}
\section[\texorpdfstring{\ml{കഴുത} (kaḻuta)}{കഴുത (kaḻuta)}]{\ml{കഴുത} (kaḻuta) — a donkey}

\label{lesson:ML-C151-kaluta}



\begin{quote}
Before the new one: say the Malayalam for garlic, then the Malayalam for a horse.

One word or a phrase: how does Malayalam say \textquotedblleft{}afternoon\textquotedblright{}? (A phrase, \emph{uccakaḻiññŭ}.)
\end{quote}

\subsection*{You'll want to know: \ml{കഴുത}}
\textbf{\ml{കഴുത}} — \emph{kaḻuta} — \textquotedblleft{}a donkey\textquotedblright{}.

\begin{grammarlens}[title={What you've built}]
One more word for this chapter.
\end{grammarlens}

\subsection*{Your turn}
\begin{itemize}
\raggedright
  \item \emph{Say it:} \emph{kaḻuta}
  \item \emph{Say it:} \emph{kaḻuta}, once more
  \item \emph{Say it:} say \emph{kaḻuta}
  \item \emph{Recall:} say the Malayalam for garlic, then the Malayalam for a horse, then say \emph{kaḻuta} again
\end{itemize}

\begin{tcolorbox}[breakable,colback=teal!4,colframe=teal!35!black,title={Before you move on}]
What does \ml{കഴുത} mean? (\textquotedblleft{}A donkey\textquotedblright{}.) Say it once more.
\end{tcolorbox}
\section[\texorpdfstring{\ml{പോത്ത്} (pōttŭ)}{പോത്ത് (pōttŭ)}]{\ml{പോത്ത്} (pōttŭ) — a buffalo}

\label{lesson:ML-C151-pottu}



\begin{quote}
Before the new one: say the Malayalam for a horse, then the Malayalam for a donkey.
\end{quote}

\subsection*{You'll want to know: \ml{പോത്ത്}}
\textbf{\ml{പോത്ത്}} — \emph{pōttŭ} — \textquotedblleft{}a buffalo\textquotedblright{}.

\begin{grammarlens}[title={What you've built}]
One more word for this chapter.
\end{grammarlens}

\subsection*{Your turn}
\begin{itemize}
\raggedright
  \item \emph{Say it:} \emph{pōttŭ}
  \item \emph{Say it:} \emph{pōttŭ}, once more
  \item \emph{Say it:} say \emph{pōttŭ}
  \item \emph{Recall:} say the Malayalam for a horse, then the Malayalam for a donkey, then say \emph{pōttŭ} again
\end{itemize}

\begin{tcolorbox}[breakable,colback=teal!4,colframe=teal!35!black,title={Before you move on}]
What does \ml{പോത്ത്} mean? (\textquotedblleft{}A buffalo\textquotedblright{}.) Say it once more.
\end{tcolorbox}
\section[\texorpdfstring{\ml{പന്നി} (panni)}{പന്നി (panni)}]{\ml{പന്നി} (panni) — a pig}

\label{lesson:ML-C151-panni}



\begin{quote}
Before the new one: say the Malayalam for a donkey, then the Malayalam for a buffalo.

\emph{Suprabhātaṁ}, \textquotedblleft{}good morning,\textquotedblright{} is the very same Sanskrit word in which neighbouring language? (Telugu.)
\end{quote}

\subsection*{You'll want to know: \ml{പന്നി}}
\textbf{\ml{പന്നി}} — \emph{panni} — \textquotedblleft{}a pig\textquotedblright{}.

\begin{grammarlens}[title={What you've built}]
One more word for this chapter.
\end{grammarlens}

\subsection*{Your turn}
\begin{itemize}
\raggedright
  \item \emph{Say it:} \emph{panni}
  \item \emph{Say it:} \emph{panni}, once more
  \item \emph{Say it:} say \emph{panni}
  \item \emph{Recall:} say the Malayalam for a donkey, then the Malayalam for a buffalo, then say \emph{panni} again
\end{itemize}

\begin{tcolorbox}[breakable,colback=teal!4,colframe=teal!35!black,title={Before you move on}]
What does \ml{പന്നി} mean? (\textquotedblleft{}A pig\textquotedblright{}.) Say it once more.
\end{tcolorbox}
\section[\texorpdfstring{\ml{സിംഹം} (simhaṁ)}{സിംഹം (simhaṁ)}]{\ml{സിംഹം} (simhaṁ) — a lion}

\label{lesson:ML-C151-simham}



\begin{quote}
Before the new one: say the Malayalam for a buffalo, then the Malayalam for a pig.
\end{quote}

\subsection*{You'll want to know: \ml{സിംഹം}}
\textbf{\ml{സിംഹം}} — \emph{simhaṁ} — \textquotedblleft{}a lion\textquotedblright{}.

\begin{grammarlens}[title={What you've built}]
One more word for this chapter.
\end{grammarlens}

\subsection*{Your turn}
\begin{itemize}
\raggedright
  \item \emph{Say it:} \emph{simhaṁ}
  \item \emph{Say it:} \emph{simhaṁ}, once more
  \item \emph{Say it:} say \emph{simhaṁ}
  \item \emph{Recall:} say the Malayalam for a buffalo, then the Malayalam for a pig, then say \emph{simhaṁ} again
\end{itemize}

\begin{tcolorbox}[breakable,colback=teal!4,colframe=teal!35!black,title={Before you move on}]
What does \ml{സിംഹം} mean? (\textquotedblleft{}A lion\textquotedblright{}.) Say it once more.
\end{tcolorbox}
